\documentclass[11pt]{article}
\usepackage[margin=1in]{geometry}
\usepackage{amsmath,amssymb}
\usepackage{microtype}
\usepackage[colorlinks=true,urlcolor=blue,citecolor=blue]{hyperref}

\title{Heat-Kernel Ces\`aro Limits and Dixmier Measurability:\\
a Later Literature Answer}
\author{Status note for arXiv:0912.2817}
\date{12 August 2026}

\begin{document}
\maketitle

\noindent\textbf{Status: literature\_already\_answered.}
This note records a separate later-paper answer; it is not an original
solution produced by the run.

\section*{Source question}

Carey, Gayral, Rennie, and Sukochev,
\emph{Integration on locally compact noncommutative spaces},
arXiv:0912.2817 \cite{CGRS}, define for \(T\in\mathcal Z_1\)
\[
 f_T(\lambda)=\lambda^{-1}
 \tau\!\left(e^{-\lambda^{-1}|T|^{-1}}\right)
\]
and use the multiplicative Ces\`aro operator
\[
 (Mh)(\lambda)=\frac1{\log\lambda}\int_1^\lambda h(t)\,\frac{dt}{t}.
\]
On source PDF page 9, after discussing independence of Dixmier traces from
the generalized limit, they state that existence of
\[
                  \lim_{\lambda\to\infty}(Mf_T)(\lambda)
 \tag{1}\label{eq:source}
\]
is sufficient for measurability and ask whether it is necessary.

\section*{Supporting answer}

Carey and Sukochev,
\emph{Measurable operators and the asymptotics of heat kernels and zeta
functions}, arXiv:1201.3661 \cite{CS}, explicitly announce that they answer
the heat-kernel measurability questions considered there.  Their Theorem 7,
on supporting PDF page 5, states that for every positive
\(A\in\mathcal M_{1,\infty}\), the following are equivalent:
\begin{enumerate}
\item \(A\) is measurable for the class of Dixmier traces;
\item the logarithmic singular-value mean has an ordinary limit;
\item
\[
 \lim_{\lambda\to\infty}
 M\!\left(\lambda\longmapsto
 \lambda^{-1}\tau(e^{-(\lambda A)^{-1}})\right)
\]
exists;
\item \(\lim_{s\to0}s\tau(A^{1+s})\) exists.
\end{enumerate}
The theorem also proves that the three limits and the common Dixmier-trace
value coincide.  Taking \(A=|T|\), condition (iii) is exactly
\eqref{eq:source}.  Thus the requested necessity holds, and in fact belongs
to a stronger necessary-and-sufficient four-way characterization.

\section*{Scope and provenance}

The source question is about the positive heat-kernel quantity attached to
\(|T|\); the supporting theorem applies to positive operators in
\(\mathcal M_{1,\infty}\), the same Marcinkiewicz ideal denoted
\(\mathcal Z_1\) in the source.  The answer paper is separate and later, and
two of its authors are authors of the source paper.  Its abstract and
introduction say that it answers the outstanding measurability/heat-kernel
questions, and Theorem 7 gives the exact equivalence.

The queue's earlier natural-question signal on the same source page asks
whether \(f_T\) may be unbounded while \(Mf_T\) is bounded.  That question is
answered immediately by the following proposition in arXiv:0912.2817 and is
therefore a same-paper extraction false positive.  This status packet records
only the distinct genuine question above.

\begin{thebibliography}{9}
\bibitem{CGRS}
A.\ L.\ Carey, V.\ Gayral, A.\ Rennie, and F.\ A.\ Sukochev,
\emph{Integration on locally compact noncommutative spaces},
J. Funct. Anal. \textbf{260} (2011), 1116--1163;
arXiv:0912.2817.

\bibitem{CS}
A.\ Carey and F.\ Sukochev,
\emph{Measurable operators and the asymptotics of heat kernels and zeta
functions},
J. Funct. Anal. \textbf{262} (2012), 4589--4599;
arXiv:1201.3661.
\end{thebibliography}

\end{document}
